\chapter*{Resumo}

A alocação de pacientes a leitos é uma das decisões mais críticas da gestão hospitalar, pois afeta a qualidade do atendimento, o conforto dos pacientes e o uso dos recursos disponíveis. Este trabalho trata o Problema de Alocação de Leitos em sua versão voltada ao planejamento tático do censo hospitalar, na qual pacientes com datas de admissão e de alta conhecidas devem ser alocados aos quartos de um hospital ao longo de um horizonte, minimizando uma soma ponderada de penalidades por especialidade inadequada, transferências e quartos ocupados por pacientes de gêneros distintos, sem exceder a capacidade dos quartos. O problema é formulado como um modelo de Programação Linear Inteira, resolvido pelo \textit{solver} Gurobi, e abordado por um Algoritmo Memético que combina um Algoritmo Genético com busca local por descida em vizinhança variável (VND) e trata a capacidade como restrição penalizada. As duas abordagens são comparadas, sob o mesmo teto de tempo, em 40 conjuntos de instâncias sintéticas de quatro portes, com 100 a 6400 pacientes. Na instância de menor porte, o \textit{solver} obtém as melhores soluções, e a heurística fica, em média, 3\% acima delas e no máximo 10\% acima do ótimo. A partir de 400 pacientes, o \textit{solver} deixa de convergir no tempo disponível, e a heurística obtém soluções com valor da função objetivo de 54\% a 95\% menor. Um limite inferior combinatório proposto neste trabalho mostra que a heurística fica a no máximo 0,8 ponto percentual do mínimo de alocações em especialidade inadequada imposto pelos dados, com coeficiente de variação inferior a 3\% entre execuções. Nas instâncias maiores, as soluções heurísticas apresentam uma violação residual de capacidade, inferior a 0,1\% da demanda de internação, que pode ser eliminada por um reparo que acrescenta no máximo 3,4\% ao custo. Os resultados indicam que a abordagem heurística é adequada ao planejamento do censo em hospitais de grande porte, nos quais o método exato não é viável no tempo disponível.

\vspace{0.5cm}
\noindent\textbf{Palavras-chave:} Alocação de leitos. Algoritmo memético. Programação linear inteira. Planejamento hospitalar. Otimização combinatória.

\chapter*{Abstract}

Assigning patients to beds is one of the most critical decisions in hospital management, as it affects the quality of care, patient comfort and the use of available resources. This work addresses the Patient-Bed Allocation problem in its version aimed at the tactical planning of the hospital census, in which patients with known admission and discharge dates must be assigned to the rooms of a hospital over a planning horizon, minimizing a weighted sum of penalties for inadequate specialty, transfers and rooms occupied by patients of different genders, without exceeding room capacity. The problem is formulated as an Integer Linear Programming model, solved with the Gurobi solver, and tackled by a Memetic Algorithm that combines a Genetic Algorithm with Variable Neighborhood Descent (VND) local search and treats capacity as a penalized constraint. Both approaches are compared, under the same time limit, on 40 sets of synthetic instances of four sizes, with 100 to 6400 patients. On the smallest instance, the solver obtains the best solutions, and the heuristic is, on average, 3\% above them and at most 10\% above the optimum. From 400 patients onward, the solver no longer converges within the available time, and the heuristic finds solutions with objective function values 54\% to 95\% lower. A combinatorial lower bound proposed in this work shows that the heuristic stays within 0.8 percentage points of the minimum number of inadequate specialty assignments imposed by the data, with a coefficient of variation below 3\% across runs. On the larger instances, the heuristic solutions show a residual capacity violation, below 0.1\% of the inpatient demand, which can be removed by a repair that adds at most 3.4\% to the cost. The results indicate that the heuristic approach is suitable for census planning in large hospitals, where the exact method is not viable within the available time.

\vspace{0.5cm}
\noindent\textbf{Keywords:} Patient-bed allocation. Memetic algorithm. Integer linear programming. Hospital planning. Combinatorial optimization.
